\documentclass[10pt,a4paper]{amsart}
\usepackage{amsmath,amssymb,amsthm,graphicx}
\usepackage{verbatim}
\newenvironment{reference}{\comment}{\endcomment}
%\newenvironment{reference}{}{}
\newenvironment{slogan}{\comment}{\endcomment}
\newenvironment{history}{\comment}{\endcomment}

% For commutative diagrams we use Xy-pic
\usepackage[all]{xy}

% We use 2cell for 2-commutative diagrams.
\xyoption{2cell}
% Theorem environments.
%
\theoremstyle{plain}
\newtheorem{theorem}[subsection]{Theorem}
\newtheorem{proposition}[subsection]{Proposition}
\newtheorem{lemma}[subsection]{Lemma}

\theoremstyle{definition}
\newtheorem{definition}[subsection]{Definition}
\newtheorem{example}[subsection]{Example}
\newtheorem{exercise}[subsection]{Exercise}
\newtheorem{situation}[subsection]{Situation}

\theoremstyle{remark}
\newtheorem{remark}[subsection]{Remark}
\newtheorem{remarks}[subsection]{Remarks}

\numberwithin{equation}{subsection}

% Macros
%
\def\lim{\mathop{\mathrm{lim}}\nolimits}
\def\colim{\mathop{\mathrm{colim}}\nolimits}
\def\Spec{\mathop{\mathrm{Spec}}}
\def\Hom{\mathop{\mathrm{Hom}}\nolimits}
\def\Ext{\mathop{\mathrm{Ext}}\nolimits}
\def\SheafHom{\mathop{\mathcal{H}\!\mathit{om}}\nolimits}
\def\SheafExt{\mathop{\mathcal{E}\!\mathit{xt}}\nolimits}
\def\Sch{\mathit{Sch}}
\def\Mor{\mathop{\mathrm{Mor}}\nolimits}
\def\Ob{\mathop{\mathrm{Ob}}\nolimits}
\def\Sh{\mathop{\mathit{Sh}}\nolimits}
\def\NL{\mathop{N\!L}\nolimits}
\def\CH{\mathop{\mathrm{CH}}\nolimits}
\def\proetale{{pro\text{-}\acute{e}tale}}
\def\etale{{\acute{e}tale}}
\def\QCoh{\mathit{QCoh}}
\def\Ker{\mathop{\mathrm{Ker}}}
\def\Im{\mathop{\mathrm{Im}}}
\def\Coker{\mathop{\mathrm{Coker}}}
\def\Coim{\mathop{\mathrm{Coim}}}

% Boxtimes
%
\DeclareMathSymbol{\boxtimes}{\mathbin}{AMSa}{"02}

%
% Macros for moduli stacks/spaces
%
\def\QCohstack{\mathcal{QC}\!\mathit{oh}}
\def\Cohstack{\mathcal{C}\!\mathit{oh}}
\def\Spacesstack{\mathcal{S}\!\mathit{paces}}
\def\Quotfunctor{\mathrm{Quot}}
\def\Hilbfunctor{\mathrm{Hilb}}
\def\Curvesstack{\mathcal{C}\!\mathit{urves}}
\def\Polarizedstack{\mathcal{P}\!\mathit{olarized}}
\def\Complexesstack{\mathcal{C}\!\mathit{omplexes}}
% \Pic is the operator that assigns to X its picard group, usage \Pic(X)
% \Picardstack_{X/B} denotes the Picard stack of X over B
% \Picardfunctor_{X/B} denotes the Picard functor of X over B
\def\Pic{\mathop{\mathrm{Pic}}\nolimits}
\def\Picardstack{\mathcal{P}\!\mathit{ic}}
\def\Picardfunctor{\mathrm{Pic}}
\def\Deformationcategory{\mathcal{D}\!\mathit{ef}}
\usepackage{cleveref}
\crefname{originalequation}{original equation}{original equations}
\Crefname{originalequation}{Original equation}{Original equations}
\makeatletter
\newlabel{equation-a-torsion}{{0AXK}{source}{}{sourcepointer-1}{}}
\newlabel{equation-a-torsion@cref}{{[originalequation][0][]0AXK}{[1][0][]source}{}{}{}}
\makeatother
\crefname{originallemma}{original lemma}{original lemmas}
\Crefname{originallemma}{Original lemma}{Original lemmas}
\makeatletter
\newlabel{fields-lemma-separable-trace-pairing}{{0BIL}{source}{}{sourcepointer-2}{}}
\newlabel{fields-lemma-separable-trace-pairing@cref}{{[originallemma][0][]0BIL}{[1][0][]source}{}{}{}}
\makeatother
\crefname{originallemma}{original lemma}{original lemmas}
\Crefname{originallemma}{Original lemma}{Original lemmas}
\makeatletter
\newlabel{flat-lemma-finite-after-blowing-up}{{0B49}{source}{}{sourcepointer-3}{}}
\newlabel{flat-lemma-finite-after-blowing-up@cref}{{[originallemma][0][]0B49}{[1][0][]source}{}{}{}}
\makeatother
\crefname{originalreference}{original reference}{original references}
\Crefname{originalreference}{Original reference}{Original references}
\makeatletter
\newlabel{lemma-bound-a-torsion}{{0AXJ}{source}{}{sourcepointer-4}{}}
\newlabel{lemma-bound-a-torsion@cref}{{[originalreference][0][]0AXJ}{[1][0][]source}{}{}{}}
\makeatother
\title{Reduction to rational singularities ascends along separable finite local extensions of normal surface singularities}
\author{The Stacks Project Authors}
\date{stacks-project:resolve.tex:local-snapshot-20261012; GNU FDL 1.2 or later}
\begin{document}
\maketitle
\noindent\textit{Editorial note.} All delivered mathematical text is the author's own, copied from the pinned Stacks Project chapter file resolve.tex; no statement or proof step was rewritten, reordered or removed.\par\medskip
\noindent\textit{Editorial note.} The counted proof is the author's complete argument: the production of a nonzero determinant from the nondegenerate trace pairing, the reduction to a normal modification finite over a modification of the source via an admissible blowup, and the comparison of the first cohomology groups through the trace map together with the bound on the remaining torsion.\par\medskip
\noindent\textit{Editorial note.} The label lemma-go-up-separable names the counted statement, and the proof environment immediately following it in the source is delivered verbatim as its proof.\par\medskip
\noindent\textit{Editorial note.} The counted statement is carried by the author's own lemma environment in the chapter source.\par\medskip
\noindent\textit{Editorial note.} Pre-open gate: python operations/agent-authoring/tools/precard\_gate.py --source resolve.tex --label lemma-go-up-separable --start 2138 --end 2202 --local-id 28084 returned RESULT=OK (L28084 free, occupancy FREE, public interval clear). Pre-assembly tag self-check: canonical tag resolved and not present in reference/public-items.jsonl.\par\medskip
\noindent\textit{Editorial note.} The statement and proof are copied from the pinned chapter file; the cited inputs are kept as external pointers with their canonical Stacks tags.\par\medskip
\noindent\textit{Allowed standard prerequisites (editorial).}\begin{itemize}
\item Nondegeneracy of the trace pairing for a finite separable extension. Fields, the lemma that the trace pairing of a finite separable field extension is nondegenerate. Supplies the nonvanishing of the determinant of the matrix of traces of a basis, the element used to bound torsion below.
\item Finiteness of a strict transform after an admissible blowup. More on Flatness, the lemma producing an admissible blowup after which the strict transform becomes finite. Supplies the reduction to a situation where the modification of the target is finite over a modification of the source, both normal.
\item Boundedness of the torsion in the first cohomology. Same chapter, the lemma bounding the torsion of the first cohomology by a fixed element. Supplies the final ingredient completing the induction of the vanishing.
\end{itemize}\medskip
\ifdefined\hypertarget\hypertarget{sourcepointer-1}{}\fi
\noindent\textit{Source reference.} Original equation 0AXK: sources/stacks-resolve/pinned-original-source.tex, source lines 1965--1965. Prerequisite pointer invoked at proof line(s) 2201; it names a standard fact of the same chapter that the counted argument uses, and every step of the counted proof built on it is retained verbatim in the delivered body.\par\medskip
\ifdefined\hypertarget\hypertarget{sourcepointer-2}{}\fi
\noindent\textit{Source reference.} Original lemma 0BIL: sources/stacks-chow/tags.txt, source lines 14494--14494. Prerequisite pointer invoked at proof line(s) 2151; it names a standard fact of a sibling chapter that the counted argument uses, and every step of the counted proof built on it is retained verbatim in the delivered body.\par\medskip
\ifdefined\hypertarget\hypertarget{sourcepointer-3}{}\fi
\noindent\textit{Source reference.} Original lemma 0B49: sources/stacks-chow/tags.txt, source lines 13992--13992. Prerequisite pointer invoked at proof line(s) 2159; it names a standard fact of a sibling chapter that the counted argument uses, and every step of the counted proof built on it is retained verbatim in the delivered body.\par\medskip
\ifdefined\hypertarget\hypertarget{sourcepointer-4}{}\fi
\noindent\textit{Source reference.} Original reference 0AXJ: sources/stacks-resolve/pinned-original-source.tex, source lines 1947--1947. Prerequisite pointer invoked at proof line(s) 2200; it names a standard fact of the same chapter that the counted argument uses, and every step of the counted proof built on it is retained verbatim in the delivered body.\par\medskip
\begin{lemma}
\label{lemma-go-up-separable}
Let $A \to B$ be a finite injective local ring map of local normal
Nagata domains of dimension $2$. Assume that the induced extension of
fraction fields is separable. If reduction to rational singularities
is possible for $A$ then it is possible for $B$.
\end{lemma}

\begin{proof}
Let $n$ be the degree of the fraction field extension $L/K$.
Let $\text{Trace}_{L/K} : L \to K$ be the trace. Since the extension is finite
separable the trace pairing $(h, g) \mapsto \text{Trace}_{L/K}(fg)$
is a nondegenerate bilinear form on $L$ over $K$. See
Fields, Lemma \ref{fields-lemma-separable-trace-pairing}.
Pick $b_1, \ldots, b_n \in B$ which form a basis of $L$ over $K$.
By the above $d = \det(\text{Trace}_{L/K}(b_ib_j)) \in A$ is nonzero.

\medskip\noindent
Let $Y \to \Spec(B)$ be a modification with $Y$ normal. We can find
a $U$-admissible blowup $X'$ of $\Spec(A)$ such that the strict transform
$Y'$ of $Y$ is finite over $X'$, see More on Flatness, Lemma
\ref{flat-lemma-finite-after-blowing-up}. Picture
$$
\xymatrix{
Y' \ar[d] \ar[r] & Y \ar[r] & \Spec(B) \ar[d] \\
X' \ar[rr] & & \Spec(A)
}
$$
After replacing $X'$ and $Y'$ by their normalizations we may assume that
$X'$ and $Y'$ are normal modifications of $\Spec(A)$ and $\Spec(B)$.
In this way we reduce to the case where there exists a commutative diagram
$$
\xymatrix{
Y \ar[d]_\pi \ar[r]_-g & \Spec(B) \ar[d] \\
X \ar[r]^-f & \Spec(A)
}
$$
with $X$ and $Y$ normal modifications of $\Spec(A)$ and $\Spec(B)$ and
$\pi$ finite.

\medskip\noindent
The trace map on $L$ over $K$ extends to a map of $\mathcal{O}_X$-modules
$\text{Trace} : \pi_*\mathcal{O}_Y \to \mathcal{O}_X$. Consider the map
$$
\Phi : \pi_*\mathcal{O}_Y \longrightarrow \mathcal{O}_X^{\oplus n},\quad
s \longmapsto (\text{Trace}(b_1s), \ldots, \text{Trace}(b_ns))
$$
This map is injective (because it is injective in the generic point)
and there is a map
$$
\mathcal{O}_X^{\oplus n} \longrightarrow \pi_*\mathcal{O}_Y,\quad
(s_1, \ldots, s_n) \longmapsto \sum b_i s_i
$$
whose composition with $\Phi$ has matrix $\text{Trace}(b_ib_j)$.
Hence the cokernel of $\Phi$ is annihilated by $d$. Thus we see that
we have an exact sequence
$$
H^0(X, \Coker(\Phi)) \to H^1(Y, \mathcal{O}_Y) \to
H^1(X, \mathcal{O}_X)^{\oplus n}
$$
Since the right hand side is bounded by assumption, it suffices to show
that the $d$-torsion in $H^1(Y, \mathcal{O}_Y)$ is bounded.
This is the content of Lemma \ref{lemma-bound-a-torsion} and
(\ref{equation-a-torsion}).
\end{proof}

\end{document}
